% 原书 p.45–p.47（PDF p.53–p.55）
% >>> p053
\section{诺特原理（Noether's Principle）}

\begin{keytakeaways}
\begin{itemize}
  \item 物理图像：对称性对应守恒量——这是物理学中最基本的原理之一；能量守恒对应
        时间平移对称性，动量守恒对应空间平移对称性，角动量守恒对应旋转对称性。
  \item 数学工具：哈密顿群作用、矩映射、Poisson 括号的斜对称性，以及
        $G$-不变哈密顿函数的定义 $H(g\cdot p)=H(p)$。
  \item 几何观点：定理 2.7.1 断言若 $H$ 是 $G$-不变的，则矩映射 $\mu$ 沿 $X_H$ 为常值，
        即 $\{\mu,H\}=0$；这是矩映射作为「守恒量包」的正式表述。
  \item 本节结论：由 $SO(3)$-不变性得到角动量守恒（例 2.7.1），由 $\T^n$-不变性得到
        $\{|z_i|^2\}$ 的守恒（例 2.7.2），从而自动给出可积系统；同时讨论 (2.34) 的
        「反向」问题：如何由守恒量反求对称群。
\end{itemize}
\end{keytakeaways}

本节我们将论证物理学中一条基本原理，它把力学系统中的对称性（例如由群作用给出的对称性）
与守恒量——随时间保持不变的量——联系起来：
\begin{equation}
  \text{对称性}\iff\text{守恒量}.
  \tag{2.34}
\end{equation}
这一原理以 E. Noether 命名。自此以后，物理学家逐渐成为群论专家，力图通过各种群的系统
研究去发现更多的守恒量。

这里我们专注于辛流形，其对称性由（哈密顿）群作用给出。如下结果说明，由这种对称性可导出
一个守恒量（它恰是定义 2.5.1 中的矩映射 $\mu$）。换言之，它刻画了 (2.34) 中「正向」的
方向。下面我们用前几节引入的记号来表述这一结果。

\begin{wthm}{2.7.1}
（Noether，参见 [168]）设 $(M,\omega)$ 是配有 Lie 群 $G$ 的哈密顿群作用的辛流形。若
哈密顿函数 $H:M\to\R$ 是 \emph{$G$-不变的}，即对任意 $g\in G$ 与 $p\in M$ 有
$H(g\cdot p)=H(p)$，则矩映射 $\mu:M\to\mathfrak{g}^*$ 沿 $X_H$ 为常值。
\end{wthm}
% [Tutor Note: 定理的结论「$\mu$ 沿 $X_H$ 为常值」是指沿 $X_H$ 的积分曲线（即哈密顿
% 流）$\mu$ 保持不变，而不是 $\mu$ 在整个 $M$ 上为常值。等价写法是
% $\mathcal{L}_{X_H}\mu=0$，或对每个 $v\in\mathfrak{g}$ 有
% $\{\mu_v,H\}=0$。注意定理只用到 $H$ 在 $G$-作用下的不变性，
% 与 $H$ 是否还依赖其他变量无关。]

% >>> p054
我们来澄清定理 2.7.1 中结论的含义：固定 $\mathfrak{g}$ 的任意一组基
$(v_1,\cdots,v_k)$，由定义 2.5.1，对 $M$ 上某个哈密顿函数 $H_{v_i}$ 有
$\langle\mu(x),v_i\rangle=H_{v_i}(x)$，于是上述结论意味着对任意 $1\le i\le k$ 有
$\dd H_{v_i}(X_H)=0$。由 Poisson 括号的斜对称性，这等价于如下关系
\begin{equation}
  \{H,H_{v_i}\}=0\qquad\text{对任意 } 1\le i\le k .
  \tag{2.35}
\end{equation}

\begin{wproof}[定理 2.7.1 的证明]
考虑曲线 $\gamma_{v_i}(t):=\exp(tX_{v_i})(p)$，它满足
\[
  \gamma_{v_i}(0)=p
  \qquad\text{以及}\qquad
  \left.\frac{\dd}{\dd t}\gamma_{v_i}(t)\right|_{t=0}=X_{v_i}(p).
\]
由于 $H$ 是 $G$-不变的，有 $H(\gamma_{v_i}(t))=H(p)$。对等式
$H(\gamma_{v_i}(t))=H(p)$ 两端求导，得
\begin{align*}
  0&=\left.\frac{\dd}{\dd t}H(\gamma_{v_i}(t))\right|_{t=0}
  =(\dd H)_{\gamma_{v_i}(0)}\left(\left.\frac{\dd}{\dd t}\gamma_{v_i}(t)\right|_{t=0}\right)\\
  &=\dd H(X_{v_i}(p)).
\end{align*}
证毕。
\end{wproof}
% [Tutor Note: 记号 $\exp(tX_{v_i})$ 在此表示\emph{向量场} $X_{v_i}$ 的时间 $t$ 流映射
% （即哈密顿流），而非 Lie 群的指数映射；它由定义 2.4.1 的 $\exp(tv_i)$ 在 $M$ 上的
% 作用给出。证明的实质是：$G$-不变性给出 $H\circ\varphi^t_{X_{v_i}}=H$，对 $t$ 求导
% 并用链式法则即得 $\dd H(X_{v_i})=0$；再由 $\iota_{X_{v_i}}\omega=-\dd H_{v_i}$
% 得 $\{\mu,H\}=0$。此处 $\gamma_{v_i}$ 的下标 $v_i$ 与 $X_{v_i}$ 对应，原文中
% $H_{v_i}$ 与 $\mu$ 的配对关系须一起核对。]

下面是与例 2.6.1 有关的定理 2.7.1 的一个直接应用。

\begin{wexam}{2.7.1}
回顾群 $SO(3)$ 通过 $A\cdot(q,p)=(Aq,Ap)$ 作用于 $(T^*\R^3,\omega_{\mathrm{can}})$，
如例 2.4.4 所示这是哈密顿作用。由例 2.5.1，其矩映射恰为角动量 $q\times p$。若哈密顿
函数 $H:T^*\R^3\to\R$（或如例 2.6.1 中 (2.31) 那样定义在 $T^*\R^3$ 的某子集上）只
依赖于范数 $|q|$ 与 $|p|$，则 $H$ 是 $SO(3)$-不变的，因为
\[
  |Aq|^2=\langle Aq,Aq\rangle=\langle q,A^{\mathrm T}Aq\rangle
  =\langle q,q\rangle=|q|^2 .
\]
于是定理 2.7.1 表明，角动量 $q\times p\in\R^3$ 沿哈密顿向量场 $X_H$ 不变（这通常表述为
角动量的守恒性质）。特别地，每个分量与 $H$ 满足 (2.35) 中的 Poisson 交换条件。这也解释
了我们在例 2.6.1 中为何猜测所需的 $H_2(=q_1p_2-q_2p_1)$（用以构成空间 Kepler 问题的
可积系统）。\hfill$//$
\end{wexam}
% [Tutor Note: 这里 $|Aq|=|q|$ 用到了 $A\in SO(3)$ 的正交性 $A^{\mathrm T}A=\mathbf{1}$；
% 同理 $|Ap|=|p|$。请注意 $SO(3)$-不变性只给出角动量的守恒（即 $\{H,q\times p\}=0$），
% 并不能单独给出 $H_2$ 与 $H$ 的交换性——后者正是 (2.35) 在 $H_{v}=A_3$ 时的特例。]

下面的例子说明，定理 2.7.1 也有助于构造可积系统。

\begin{wexam}{2.7.2}
回顾 $\T^n$ 通过
\[
  (t_1,\cdots,t_n)\cdot(z_1,\cdots,z_n)
  =\big(\ee^{\sqrt{-1}t_1}z_1,\cdots,\ee^{\sqrt{-1}t_n}z_n\big),
  \qquad t_i\in[0,2\pi]
\]
作用于 $(\C^n,\omega_{\mathrm{std}})$。对任意只依赖于 $|z_i|$ 的哈密顿函数
$H:\C^n\to\R$，它是 $\T^n$-不变的，因为 $|z_i|=|\ee^{\sqrt{-1}t_i}z_i|$。于是定理 2.7.1
表明 $\mu$ 沿 $X_H$ 为常值。更确切地说，由例 2.5.2 知
$\mu(z_1,\cdots,z_n)=\frac12(|z_1|^2,\cdots,|z_n|^2)$，故有
\[
  \{H,|z_i|^2\}=0\qquad\text{对 } i=1,\cdots,n .
\]
% >>> p055
由于 Poisson 括号 $\{|z_i|^2,|z_j|^2\}=0$ 对任意 $1\le i,j\le n$ 成立，可知这样的
哈密顿系统 $(\C^n,\omega_{\mathrm{std}},H)$ 总是可积的。\hfill$//$
\end{wexam}
% [Tutor Note: 这里「总是可积」依赖于 $H$ 只沿 $|z_i|$（等价地 $|z_i|^2$）变化这一
% 假设：此时 $\{\dd|z_i|^2\}$ 在 $z_i\neq0$ 的开稠密集上线性无关，且两两 Poisson 交换，
% 加上 $\{H,|z_i|^2\}=0$ 即满足定义 2.6.1。但若 $H$ 还依赖 $z_i$ 的辐角，则不变性
% 与交换性均可能失效。另注：$\mu$ 在此取的是 $|z_i|^2/2$ 而非 $|z_i|^2$，二者只差常数因子，
% 不影响括号为零的结论。]

若把 (2.34) 中 Noether 原理的「反向」读法用于此，则从任一守恒量出发，应当能求得某种
对称性（由群作用给出）。最著名的守恒量是时间流动时的总能量，其对应的群作用由 $\R$
给出，即时间变量上的平移。与这一简单情形形成对照，回顾我们在例 2.6.1 中引入的
Runge--Lenz 向量 $V_{\mathrm{RL}}$。由练习 2.6.1 知 $V_{\mathrm{RL}}$ 的每个分量都是
守恒量，而探索与 $V_{\mathrm{RL}}$ 相对应的对称性或群作用是很有意思的。
% [Tutor Note: (2.34) 的「反向」问题（Noether 逆问题）一般没有唯一答案：给定守恒量，
% 与之对应的对称群可能并不显式，甚至（在可积系统情形）只能由作用量–角坐标隐含给出。
% 对 Kepler 问题，$H$、$A$、$V_{\mathrm{RL}}$ 生成的对称群是 $SO(4)$（束缚态，负能量）或
% 与其相关的非紧群（散射态）——这是「隐藏对称性」的经典例子，可参见
% Arnol'd《经典力学的数学方法》有关 Kepler 问题的章节。]
